
\section{12世紀の熊野寮について}
\rightline{（文責：水城秀宗）}
　熊野寮は2025年をもって建寮60周年を迎え、2026年には61年、人間でいうところの還暦となる。吉田寮現棟は1913年に竣工され2026年では113年、吉田寮食堂はさらに古く137年の歴史をもつことになる。熊野寮東側の川端警察署は103年、南側の夷川発電所は112年、北側の京都丸太町教会は122年、西側のスイミングスクール踏水会は130年。さらに遡って近世の絵地図を見れば、聖護院や熊野神社を容易に見つけることができる。そもそも京都は千年以上にわたり都市・政庁として歴史の舞台となっており、至るところの記念碑や史跡がそれを示している。しかし、現在の熊野寮にあたるまさにその場所の歴史について、どれほどの人が知っているだろうか。本稿では、熊野寮所在地における特筆すべき事柄として、中世初期の様相について解説する。
　元永元年（1118）、白河上皇が政務を執り行う院御所として白河北殿が建設された。「北殿」の呼称は南に隣接する院御所・白河泉殿（南殿）と対になるものである。白河北殿は大治四年（1129）拡張、天養元年（1144）焼失のち再建との経歴をたどりつつ、白河・鳥羽・崇徳の三上皇が政務を行うにあたって用いられている。上皇が政務を執り行う政治体制は院政とよばれ、院近臣の伸長、知行国制の確立、中世荘園制の展開、統治思想の変容などの帰結により日本中世を大きく規定した。その院政の舞台のひとつとなった白河北殿が現在の熊野寮に所在したことが、熊野寮敷地内北西端に立つ「白河北殿跡」の石碑に銘記されている。熊野寮近辺は院政の一大拠点であり、皇族によって院御所や御願寺が多く創建された。吉田寮近辺には鳥羽上皇后泰子の御願寺・福勝院が所在したと推定される。川端警察署近辺には鳥羽上皇の御願寺・得長寿院が所在した。夷川発電所は白河南殿、丸太町教会・踏水会は白河北殿の敷地に接している。岡崎一帯に建立された六つの御願寺（総称：六勝寺）は国家による仏教の再編に大きな役割を果たしている。また、白河上皇は修験道・熊野信仰に帰依し、聖護院・熊野神社の建立に参与している。これらの建設に協力し地位を高めたのが、平忠盛ら院近臣であった。
　さらに白河北殿は、日本史上でも重要な合戦の舞台となっている。崇徳上皇・左大臣藤原頼長らが後白河天皇・藤原忠通らと対立を深めた末に勃発した保元の乱である。保元元年（1156）七月十日、崇徳上皇らは白河北殿に籠って戦備を整え、源為義・平家弘をはじめとする武者が参陣したことが『兵範記』に記録されている。対する御白河天皇陣営は十一日早朝、白河北殿を襲撃した。忠盛の子である平清盛、頼朝・義経らの父である源義朝、鵺退治の逸話で知られる源頼政などが攻撃に参加している。周知の通り御白河天皇らは勝利を納め、敗北した崇徳上皇は讃岐に流され、頼長は敗走中に死亡し、家弘・為義らは処刑の憂き目に遭う。一方で、勝者として生き残った者たちの多くはその後の動乱にそれぞれ活躍を見せている。平治元年（1159）平治の乱では平清盛が源義朝を打倒し、治承四年（1180）には源頼政が平氏討伐を図って挙兵、続いて頼朝らの挙兵から平氏滅亡・鎌倉幕府成立へとつながる。保元の乱を題材にした軍記『保元物語』では、源義朝に付き従う者として上総広常や千葉常胤ら鎌倉幕府創設に関わった人々が挙げられるなど、一連の戦乱は有機的に接続するものとして広く認識されている。
　保元の乱は『愚管抄』などにおいて「ムサ（武者）ノ世」の到来の契機として述べられ、のちの歴史観にも大きな影響を与えている。『保元物語』に述べられる「皇を取って民となし、民を皇となさん」という崇徳院（崇徳上皇）の呪詛、『明月記』にみられる崇徳院怨霊と鎌倉とを関連づける風聞、『関東御式目』奥書の北条泰時を「崇徳院の後身」とみなす言説など、崇徳院と朝廷の衰微・幕府の興隆を紐付ける思想が鎌倉時代以降さまざまな史料に残っている。明治維新期には孝明天皇・明治天皇の命により、堀川今出川に白峯神宮が創建され崇徳院が祀られている。幕政の終焉・王政への復古に際しても崇徳院怨霊の存在が意識されているのである。天皇親政を理想とする理論的枠組みのなかでは、武家政権である幕府が朝廷に比肩する権威をもつ事態は政体上の転倒であり、その原因のひとつとして保元の乱に敗北した崇徳院の怨霊が配置された。また、崇徳上皇方に参戦した源為朝も、数多くの伝承にその名を残している。『保元物語』での合戦描写は大部分が彼を中心に構成されており、敗戦後は伊豆大島に配流されたと述べられる。伊豆諸島や琉球・奄美群島においては為朝渡来の伝承が多く残り、近世には為朝を琉球王朝の祖として語るものもある。『椿説弓張月』はこれらを題材にした読本として著名である。これらの為朝伝承は、日本の空間構成における終端としての島嶼にかかるものであり、周辺地域への侵攻による支配領域の遷移にともなって通時代的に転変を遂げている。
　以上の通り現熊野寮は、中世初期の政務の拠点のひとつであり、歴史・伝承に名を残す人々が交錯した戦争史跡であり、歴史観・世界観における大きな転換の起点であったといえる。贔屓目かつ牽強付会な見解ではあるが、僅かにでも熊野寮に対する知見を深めることに資することができれば幸いである。

\begin{tcolorbox}[colback=white,colframe=black,title=参考文献・引用文献]
河内祥輔 『保元の乱・平治の乱』吉川弘文館、2002\\
佐藤雄基「鎌倉時代における天皇像と将軍・得宗」『史学雑誌』129、2020\\
関幸彦『英雄伝説の日本史』講談社、2019\\
樋口健太郎「院政の成立」『日本史の現在3 中世』山川出版社、2024\\
美川圭『院政：もうひとつの天皇制』中央公論新社、2016\\
「【白河北殿跡】京都市：左京区／聖護院村」『日本歴史地名大系』平凡社、1981\\
『新日本古典文学大系　保元物語』岩波書店、1961
\end{tcolorbox}